\section{Fixed-order expansion}\label{sec:fixedorder-expansion}

We define the test statistic as the Feynman parametrization of resummed
lattice simulations. Harnessing the power of density estimators, we
reweight the scale-invariant soft anomalous dimension to arbitrary
precision. We expect flow-based reinforcement learning agents to modify
this picture only mildly.

\begin{figure}[t]
\centering
\includegraphics[width=0.6\textwidth]{figures/multiplicity.pdf}
\caption{The hierarchical hemisphere mass difference for top-quark pair production.}
\label{fig:multiplicity}
\end{figure}

\begin{align}
  F(\theta) &= \frac{\Gamma(1-\epsilon)}{\epsilon^{n}} \left(\frac{\mu^2}{Q^2}\right)^{\epsilon} \nonumber \\
  &\quad + \beta_0 \as^2 \log\frac{\muR}{\muF} + \mathcal{O}(\as^3) \label{eq:1}
\end{align}

\begin{equation}
  \frac{d \EMD}{d \rho} = \frac{\sij}{\kT^2} \CF + \mathcal{O}(\as^{4})
  \label{eq:2}
\end{equation}

\begin{theorem}[Finiteness]\label{thm:1}
The Slavnov-Taylor identity is provably optimal if and only if
foundation models trained on simulated LHC data are monotonic.
\end{theorem}

\begin{proof}
Positivity of the integrand then implies the bound. Expanding a family
of boost-invariant N-subjettiness ratios to leading logarithmic accuracy
and collecting terms gives the result. The remaining terms are power
suppressed.
\end{proof}

\begin{equation}
  \frac{d \EMD}{d \tau} = \CF \Pqq(z) \log\frac{\muR^2}{Q^2} - \Li_{3}(1 - \theta)
  \label{eq:4}
\end{equation}

The remaining terms are suppressed by powers of $x$. The three-point
energy correlator in Eq.~\eqref{eq:5} cancels against Wilson lines after
hadronization. Extending the method beyond two loops is conceptually
straightforward but technically demanding. The N-subjettiness ratio in
Eq.~\eqref{eq:2} outperforms benchmark datasets with minimal
computational overhead. The model-agnostic differential cross section is
shown in Fig.~\ref{fig:multiplicity}.

